% Seminar 4 - Page 4. Fragment for
% % Seminar 4 - Page 4. Fragment for
% % Seminar 4 - Page 4. Fragment for
% % Seminar 4 - Page 4. Fragment for
% \input{seminars/pages/seminar4-pagina4-en}; no preamble or forced page break.
% EDITORIAL NOTES:
% - The opening "Show that" statement appears to compare multiplicative
%   monoids mN+1 and tN+1 with groups U(Z_m), U(Z_t), presumably the
%   groups of units of the residue-class rings. The handwritten symbol
%   resembles a large U; this identification is provisional. No proof
%   of this opening statement is present on the manuscript page.
% - Exercise 9, AS WRITTEN, requires a left identity ex=x and a RIGHT
%   inverse xx'=e for each x. That combination does not characterise
%   groups (a right-zero semigroup provides a counterexample).
%   The handwritten PROOF instead chooses LEFT inverses x'x=e and
%   x''x'=e, then derives xx'=e and xe=x. We retain the two distinct
%   readings faithfully rather than silently repairing the theorem.
% - Exercise 10 claims a finite subset H of a group is a subgroup iff
%   it is closed under products; one must also require H non-empty.
%   The proof implicitly uses h in H and hence this condition.
% - The number "1" and "e" both denote the group identity in the notes.
% - Scribbled-out marks, highlighting strokes and faint reverse-page
%   impressions have not been reconstructed.

\subsection*{Seminar 4 - Page 4}

\noindent\underline{Show that:}
\[
\begin{gathered}
 (m\mathbb N+1,\cdot)\cong(t\mathbb N+1,\cdot),
 \qquad m,t\in\mathbb N \\[2pt]
 \Longleftrightarrow\\[-2pt]
 (\mathcal U(\mathbb Z_m),\cdot)
   \cong (\mathcal U(\mathbb Z_t),\cdot).
\end{gathered}
\]
\noindent (The manuscript labels $\mathcal U(\mathbb Z_m)$ as a group.)

\medskip
\noindent\underline{Exercise 9.}
Let $(G,\cdot)$ be a semigroup. Then $G$ is a group if and only if
\begin{enumerate}
 \item there exists $e\in G$ such that
       $ex=x$ for all $x\in G$;
 \item for every $x\in G$ there exists $x'\in G$ with
       $xx'=e$ (as written in the statement).
\end{enumerate}

\noindent\underline{Proof (converse, as written in the notes).}
Let $x\in G$. Choose $x'\in G$ such that $x'x=e$, and
choose $x''\in G$ such that $x''x'=e$. Then
\[
\begin{aligned}
 xx'&=(ex)x'=(x''x')xx'
       =x''(x'x)x'=x''ex'=x''x'=e,\\
 xe&=x(x'x)=(xx')x=ex=x.
\end{aligned}
\]

\medskip
\noindent\underline{Exercise 10.}
Let $(G,\cdot)$ be a group and $H\subseteq G$ a finite set. Then
\[
 H\leq G\quad\Longleftrightarrow\quad
 H\text{ is closed under multiplication}.
\]

\noindent\underline{Proof.}
The forward implication is clear. Conversely, let $h\in H$ and put
\[
 hH=\{hx\mid x\in H\},\qquad
 H^2=\{xy\mid x,y\in H\}.
\]
Since $H$ is closed under multiplication,
\[
 hH\subseteq H^2\subseteq H.
\]
The map
\[
 \varphi_h:H\longrightarrow hH,\qquad
 \varphi_h(x)=hx
\]
is bijective. Since $H$ is finite, $|hH|=|H|$; hence
\[
 hH=H^2=H.
\]
As $h\in H=hH$, there exists $x\in H$ such that
$h=hx$, so $x=1_G$. Thus $1_G\in H$. Furthermore,
\[
 1_G\in hH\ \Longrightarrow\
  1_G=hx\ (x\in H)\ \Longrightarrow\
  x=h^{-1}\in H.
\]

\medskip
\noindent\underline{Cantor-Bernstein theorem.}
If $\alpha=|A|$, $\beta=|B|$, and
\[
 \alpha\leq\beta,\qquad\beta\leq\alpha,
 \qquad\text{then}\qquad \alpha=\beta.
\]
; no preamble or forced page break.
% EDITORIAL NOTES:
% - The opening "Show that" statement appears to compare multiplicative
%   monoids mN+1 and tN+1 with groups U(Z_m), U(Z_t), presumably the
%   groups of units of the residue-class rings. The handwritten symbol
%   resembles a large U; this identification is provisional. No proof
%   of this opening statement is present on the manuscript page.
% - Exercise 9, AS WRITTEN, requires a left identity ex=x and a RIGHT
%   inverse xx'=e for each x. That combination does not characterise
%   groups (a right-zero semigroup provides a counterexample).
%   The handwritten PROOF instead chooses LEFT inverses x'x=e and
%   x''x'=e, then derives xx'=e and xe=x. We retain the two distinct
%   readings faithfully rather than silently repairing the theorem.
% - Exercise 10 claims a finite subset H of a group is a subgroup iff
%   it is closed under products; one must also require H non-empty.
%   The proof implicitly uses h in H and hence this condition.
% - The number "1" and "e" both denote the group identity in the notes.
% - Scribbled-out marks, highlighting strokes and faint reverse-page
%   impressions have not been reconstructed.

\subsection*{Seminar 4 - Page 4}

\noindent\underline{Show that:}
\[
\begin{gathered}
 (m\mathbb N+1,\cdot)\cong(t\mathbb N+1,\cdot),
 \qquad m,t\in\mathbb N \\[2pt]
 \Longleftrightarrow\\[-2pt]
 (\mathcal U(\mathbb Z_m),\cdot)
   \cong (\mathcal U(\mathbb Z_t),\cdot).
\end{gathered}
\]
\noindent (The manuscript labels $\mathcal U(\mathbb Z_m)$ as a group.)

\medskip
\noindent\underline{Exercise 9.}
Let $(G,\cdot)$ be a semigroup. Then $G$ is a group if and only if
\begin{enumerate}
 \item there exists $e\in G$ such that
       $ex=x$ for all $x\in G$;
 \item for every $x\in G$ there exists $x'\in G$ with
       $xx'=e$ (as written in the statement).
\end{enumerate}

\noindent\underline{Proof (converse, as written in the notes).}
Let $x\in G$. Choose $x'\in G$ such that $x'x=e$, and
choose $x''\in G$ such that $x''x'=e$. Then
\[
\begin{aligned}
 xx'&=(ex)x'=(x''x')xx'
       =x''(x'x)x'=x''ex'=x''x'=e,\\
 xe&=x(x'x)=(xx')x=ex=x.
\end{aligned}
\]

\medskip
\noindent\underline{Exercise 10.}
Let $(G,\cdot)$ be a group and $H\subseteq G$ a finite set. Then
\[
 H\leq G\quad\Longleftrightarrow\quad
 H\text{ is closed under multiplication}.
\]

\noindent\underline{Proof.}
The forward implication is clear. Conversely, let $h\in H$ and put
\[
 hH=\{hx\mid x\in H\},\qquad
 H^2=\{xy\mid x,y\in H\}.
\]
Since $H$ is closed under multiplication,
\[
 hH\subseteq H^2\subseteq H.
\]
The map
\[
 \varphi_h:H\longrightarrow hH,\qquad
 \varphi_h(x)=hx
\]
is bijective. Since $H$ is finite, $|hH|=|H|$; hence
\[
 hH=H^2=H.
\]
As $h\in H=hH$, there exists $x\in H$ such that
$h=hx$, so $x=1_G$. Thus $1_G\in H$. Furthermore,
\[
 1_G\in hH\ \Longrightarrow\
  1_G=hx\ (x\in H)\ \Longrightarrow\
  x=h^{-1}\in H.
\]

\medskip
\noindent\underline{Cantor-Bernstein theorem.}
If $\alpha=|A|$, $\beta=|B|$, and
\[
 \alpha\leq\beta,\qquad\beta\leq\alpha,
 \qquad\text{then}\qquad \alpha=\beta.
\]
; no preamble or forced page break.
% EDITORIAL NOTES:
% - The opening "Show that" statement appears to compare multiplicative
%   monoids mN+1 and tN+1 with groups U(Z_m), U(Z_t), presumably the
%   groups of units of the residue-class rings. The handwritten symbol
%   resembles a large U; this identification is provisional. No proof
%   of this opening statement is present on the manuscript page.
% - Exercise 9, AS WRITTEN, requires a left identity ex=x and a RIGHT
%   inverse xx'=e for each x. That combination does not characterise
%   groups (a right-zero semigroup provides a counterexample).
%   The handwritten PROOF instead chooses LEFT inverses x'x=e and
%   x''x'=e, then derives xx'=e and xe=x. We retain the two distinct
%   readings faithfully rather than silently repairing the theorem.
% - Exercise 10 claims a finite subset H of a group is a subgroup iff
%   it is closed under products; one must also require H non-empty.
%   The proof implicitly uses h in H and hence this condition.
% - The number "1" and "e" both denote the group identity in the notes.
% - Scribbled-out marks, highlighting strokes and faint reverse-page
%   impressions have not been reconstructed.

\subsection*{Seminar 4 - Page 4}

\noindent\underline{Show that:}
\[
\begin{gathered}
 (m\mathbb N+1,\cdot)\cong(t\mathbb N+1,\cdot),
 \qquad m,t\in\mathbb N \\[2pt]
 \Longleftrightarrow\\[-2pt]
 (\mathcal U(\mathbb Z_m),\cdot)
   \cong (\mathcal U(\mathbb Z_t),\cdot).
\end{gathered}
\]
\noindent (The manuscript labels $\mathcal U(\mathbb Z_m)$ as a group.)

\medskip
\noindent\underline{Exercise 9.}
Let $(G,\cdot)$ be a semigroup. Then $G$ is a group if and only if
\begin{enumerate}
 \item there exists $e\in G$ such that
       $ex=x$ for all $x\in G$;
 \item for every $x\in G$ there exists $x'\in G$ with
       $xx'=e$ (as written in the statement).
\end{enumerate}

\noindent\underline{Proof (converse, as written in the notes).}
Let $x\in G$. Choose $x'\in G$ such that $x'x=e$, and
choose $x''\in G$ such that $x''x'=e$. Then
\[
\begin{aligned}
 xx'&=(ex)x'=(x''x')xx'
       =x''(x'x)x'=x''ex'=x''x'=e,\\
 xe&=x(x'x)=(xx')x=ex=x.
\end{aligned}
\]

\medskip
\noindent\underline{Exercise 10.}
Let $(G,\cdot)$ be a group and $H\subseteq G$ a finite set. Then
\[
 H\leq G\quad\Longleftrightarrow\quad
 H\text{ is closed under multiplication}.
\]

\noindent\underline{Proof.}
The forward implication is clear. Conversely, let $h\in H$ and put
\[
 hH=\{hx\mid x\in H\},\qquad
 H^2=\{xy\mid x,y\in H\}.
\]
Since $H$ is closed under multiplication,
\[
 hH\subseteq H^2\subseteq H.
\]
The map
\[
 \varphi_h:H\longrightarrow hH,\qquad
 \varphi_h(x)=hx
\]
is bijective. Since $H$ is finite, $|hH|=|H|$; hence
\[
 hH=H^2=H.
\]
As $h\in H=hH$, there exists $x\in H$ such that
$h=hx$, so $x=1_G$. Thus $1_G\in H$. Furthermore,
\[
 1_G\in hH\ \Longrightarrow\
  1_G=hx\ (x\in H)\ \Longrightarrow\
  x=h^{-1}\in H.
\]

\medskip
\noindent\underline{Cantor-Bernstein theorem.}
If $\alpha=|A|$, $\beta=|B|$, and
\[
 \alpha\leq\beta,\qquad\beta\leq\alpha,
 \qquad\text{then}\qquad \alpha=\beta.
\]
; no preamble or forced page break.
% EDITORIAL NOTES:
% - The opening "Show that" statement appears to compare multiplicative
%   monoids mN+1 and tN+1 with groups U(Z_m), U(Z_t), presumably the
%   groups of units of the residue-class rings. The handwritten symbol
%   resembles a large U; this identification is provisional. No proof
%   of this opening statement is present on the manuscript page.
% - Exercise 9, AS WRITTEN, requires a left identity ex=x and a RIGHT
%   inverse xx'=e for each x. That combination does not characterise
%   groups (a right-zero semigroup provides a counterexample).
%   The handwritten PROOF instead chooses LEFT inverses x'x=e and
%   x''x'=e, then derives xx'=e and xe=x. We retain the two distinct
%   readings faithfully rather than silently repairing the theorem.
% - Exercise 10 claims a finite subset H of a group is a subgroup iff
%   it is closed under products; one must also require H non-empty.
%   The proof implicitly uses h in H and hence this condition.
% - The number "1" and "e" both denote the group identity in the notes.
% - Scribbled-out marks, highlighting strokes and faint reverse-page
%   impressions have not been reconstructed.

\subsection*{Seminar 4 - Page 4}

\noindent\underline{Show that:}
\[
\begin{gathered}
 (m\mathbb N+1,\cdot)\cong(t\mathbb N+1,\cdot),
 \qquad m,t\in\mathbb N \\[2pt]
 \Longleftrightarrow\\[-2pt]
 (\mathcal U(\mathbb Z_m),\cdot)
   \cong (\mathcal U(\mathbb Z_t),\cdot).
\end{gathered}
\]
\noindent (The manuscript labels $\mathcal U(\mathbb Z_m)$ as a group.)

\medskip
\noindent\underline{Exercise 9.}
Let $(G,\cdot)$ be a semigroup. Then $G$ is a group if and only if
\begin{enumerate}
 \item there exists $e\in G$ such that
       $ex=x$ for all $x\in G$;
 \item for every $x\in G$ there exists $x'\in G$ with
       $xx'=e$ (as written in the statement).
\end{enumerate}

\noindent\underline{Proof (converse, as written in the notes).}
Let $x\in G$. Choose $x'\in G$ such that $x'x=e$, and
choose $x''\in G$ such that $x''x'=e$. Then
\[
\begin{aligned}
 xx'&=(ex)x'=(x''x')xx'
       =x''(x'x)x'=x''ex'=x''x'=e,\\
 xe&=x(x'x)=(xx')x=ex=x.
\end{aligned}
\]

\medskip
\noindent\underline{Exercise 10.}
Let $(G,\cdot)$ be a group and $H\subseteq G$ a finite set. Then
\[
 H\leq G\quad\Longleftrightarrow\quad
 H\text{ is closed under multiplication}.
\]

\noindent\underline{Proof.}
The forward implication is clear. Conversely, let $h\in H$ and put
\[
 hH=\{hx\mid x\in H\},\qquad
 H^2=\{xy\mid x,y\in H\}.
\]
Since $H$ is closed under multiplication,
\[
 hH\subseteq H^2\subseteq H.
\]
The map
\[
 \varphi_h:H\longrightarrow hH,\qquad
 \varphi_h(x)=hx
\]
is bijective. Since $H$ is finite, $|hH|=|H|$; hence
\[
 hH=H^2=H.
\]
As $h\in H=hH$, there exists $x\in H$ such that
$h=hx$, so $x=1_G$. Thus $1_G\in H$. Furthermore,
\[
 1_G\in hH\ \Longrightarrow\
  1_G=hx\ (x\in H)\ \Longrightarrow\
  x=h^{-1}\in H.
\]

\medskip
\noindent\underline{Cantor-Bernstein theorem.}
If $\alpha=|A|$, $\beta=|B|$, and
\[
 \alpha\leq\beta,\qquad\beta\leq\alpha,
 \qquad\text{then}\qquad \alpha=\beta.
\]
